\documentclass[UTF8,a4paper,12pt]{ctexart}

\usepackage{geometry}
\usepackage{xcolor}
\usepackage{listings}
\usepackage{hyperref}
\usepackage{longtable}
\usepackage{array}
\usepackage{enumitem}
\usepackage{amsmath}
\usepackage{booktabs}

\geometry{left=2.2cm,right=2.2cm,top=2.2cm,bottom=2.2cm}

\hypersetup{
    colorlinks=true,
    linkcolor=blue,
    urlcolor=blue
}
\lstdefinelanguage{CMake}{
    morekeywords={
        cmake_minimum_required,
        project,
        find_package,
        add_executable,
        ament_target_dependencies,
        install,
        ament_package,
        TARGETS,
        DESTINATION,
        REQUIRED,
        VERSION
    },
    sensitive=true,
    morecomment=[l]{\#},
    morestring=[b]"
}

\definecolor{codebg}{RGB}{248,248,248}
\definecolor{keywordcolor}{RGB}{0,0,150}
\definecolor{commentcolor}{RGB}{0,120,0}
\definecolor{stringcolor}{RGB}{160,40,40}

\lstset{
    backgroundcolor=\color{codebg},
    basicstyle=\ttfamily\small,
    keywordstyle=\color{keywordcolor}\bfseries,
    commentstyle=\color{commentcolor},
    stringstyle=\color{stringcolor},
    numbers=left,
    numberstyle=\tiny,
    stepnumber=1,
    numbersep=8pt,
    frame=single,
    breaklines=true,
    tabsize=4,
    showstringspaces=false,
    columns=fullflexible
}

\title{\Huge Linux 与 ROS 2 嵌入式期末复习资料}
\author{复习讲义}
\date{\today}

\begin{document}

\maketitle
\tableofcontents
\newpage

\section{Linux 基础概念}

\subsection{什么是 Linux}

Linux 是一种类 Unix 操作系统内核，常见发行版包括 Ubuntu、Debian、CentOS、Fedora、Arch Linux 等。

Linux 常用于：
\begin{itemize}
    \item 服务器系统；
    \item 嵌入式系统；
    \item 机器人系统；
    \item 云计算平台；
    \item 开发环境。
\end{itemize}

Linux 系统一般由以下部分组成：
\begin{itemize}
    \item Linux Kernel：内核，负责进程、内存、文件系统、设备驱动、网络等；
    \item Shell：命令解释器，例如 bash、zsh；
    \item GNU 工具：如 ls、cp、gcc、make；
    \item 文件系统：组织文件和设备；
    \item 用户空间程序：用户运行的程序。
\end{itemize}

\subsection{Linux 的核心思想}

\begin{enumerate}
    \item 一切皆文件。
    \item 小工具组合完成复杂任务。
    \item 使用文本文件作为配置文件。
    \item 多用户、多任务。
    \item 权限控制严格。
\end{enumerate}

\subsection{Linux 目录结构}

\begin{longtable}{p{4cm}p{10cm}}
\toprule
目录 & 作用 \\
\midrule
/ & 根目录，所有目录的起点 \\
/bin & 基本命令，如 ls、cp、mv \\
/sbin & 系统管理命令，如 ifconfig、reboot \\
/etc & 配置文件目录 \\
/home & 普通用户家目录 \\
/root & root 用户家目录 \\
/usr & 用户程序和库文件 \\
/var & 日志、缓存、可变数据 \\
/tmp & 临时文件 \\
/dev & 设备文件，如 /dev/ttyUSB0 \\
/proc & 虚拟文件系统，内核和进程信息 \\
/sys & 虚拟文件系统，设备和驱动信息 \\
/boot & 内核、启动文件 \\
/lib & 系统库文件 \\
/mnt & 临时挂载点 \\
/media & 自动挂载设备，如 U 盘 \\
\bottomrule
\end{longtable}

嵌入式中常见的重要目录：
\begin{itemize}
    \item \texttt{/dev}：串口、I2C、SPI、GPIO 等设备文件；
    \item \texttt{/proc}：查看 CPU、内存、进程等信息；
    \item \texttt{/sys}：访问设备、驱动、GPIO、PWM 等接口；
    \item \texttt{/etc/init.d} 或 \texttt{systemd}：启动脚本和服务管理。
\end{itemize}

\section{Linux 常用命令}

\subsection{命令基本格式}

Linux 命令一般格式：

\begin{lstlisting}[language=bash]
command [options] [arguments]
\end{lstlisting}

例如：

\begin{lstlisting}[language=bash]
ls -l /home
\end{lstlisting}

含义：
\begin{itemize}
    \item \texttt{ls}：命令；
    \item \texttt{-l}：选项，表示长格式显示；
    \item \texttt{/home}：参数，表示查看 /home 目录。
\end{itemize}

\subsection{获取帮助}

\begin{lstlisting}[language=bash]
man ls
ls --help
info ls
\end{lstlisting}

例如查看 \texttt{cp} 命令帮助：

\begin{lstlisting}[language=bash]
man cp
cp --help
\end{lstlisting}

\subsection{目录操作}

\begin{longtable}{p{4cm}p{10cm}}
\toprule
命令 & 作用 \\
\midrule
pwd & 显示当前所在目录 \\
cd & 切换目录 \\
ls & 查看目录内容 \\
mkdir & 创建目录 \\
rmdir & 删除空目录 \\
tree & 以树形显示目录结构 \\
\bottomrule
\end{longtable}

示例：

\begin{lstlisting}[language=bash]
pwd
cd /home
cd ..
cd ~
mkdir test
mkdir -p dir1/dir2/dir3
rmdir test
ls
ls -l
ls -a
ls -lh
\end{lstlisting}

解释：
\begin{itemize}
    \item \texttt{cd ..}：返回上一级目录；
    \item \texttt{cd \textasciitilde}：进入当前用户家目录；
    \item \texttt{mkdir -p}：递归创建多级目录；
    \item \texttt{ls -a}：显示隐藏文件；
    \item \texttt{ls -lh}：以人类可读方式显示文件大小。
\end{itemize}

\subsection{文件操作}

\begin{longtable}{p{4cm}p{10cm}}
\toprule
命令 & 作用 \\
\midrule
touch & 创建空文件或更新时间戳 \\
cp & 复制文件或目录 \\
mv & 移动或重命名文件 \\
rm & 删除文件或目录 \\
cat & 查看文件内容 \\
less & 分页查看文件 \\
head & 查看文件开头 \\
tail & 查看文件结尾 \\
file & 查看文件类型 \\
stat & 查看文件详细属性 \\
\bottomrule
\end{longtable}

示例：

\begin{lstlisting}[language=bash]
touch a.txt
cp a.txt b.txt
cp -r dir1 dir2
mv b.txt c.txt
rm c.txt
rm -r dir2
rm -rf dir2
cat a.txt
less a.txt
head -n 5 a.txt
tail -n 10 a.txt
tail -f /var/log/syslog
file a.txt
stat a.txt
\end{lstlisting}

注意：
\begin{itemize}
    \item \texttt{rm -rf} 非常危险，表示强制递归删除；
    \item \texttt{tail -f} 常用于实时查看日志；
    \item 嵌入式调试时常用 \texttt{dmesg} 和 \texttt{tail -f} 查看驱动或系统日志。
\end{itemize}

\section{文件查看与文本处理}

\subsection{grep 查找文本}

\texttt{grep} 用于在文件中查找匹配内容。

\begin{lstlisting}[language=bash]
grep "main" test.c
grep -n "error" log.txt
grep -i "warning" log.txt
grep -r "TODO" ./src
grep -v "debug" log.txt
\end{lstlisting}

常用选项：
\begin{itemize}
    \item \texttt{-n}：显示行号；
    \item \texttt{-i}：忽略大小写；
    \item \texttt{-r}：递归查找；
    \item \texttt{-v}：反向匹配。
\end{itemize}

例子：查找系统日志中的错误：

\begin{lstlisting}[language=bash]
grep -i "error" /var/log/syslog
\end{lstlisting}

\subsection{find 查找文件}

\begin{lstlisting}[language=bash]
find /home -name "*.c"
find . -type f -name "*.txt"
find . -type d -name "build"
find . -size +10M
find . -mtime -7
\end{lstlisting}

含义：
\begin{itemize}
    \item \texttt{-name}：按文件名查找；
    \item \texttt{-type f}：普通文件；
    \item \texttt{-type d}：目录；
    \item \texttt{-size +10M}：大于 10 MB 的文件；
    \item \texttt{-mtime -7}：7 天内修改过的文件。
\end{itemize}

删除所有 build 目录：

\begin{lstlisting}[language=bash]
find . -type d -name "build" -exec rm -rf {} \;
\end{lstlisting}

\subsection{sed 文本替换}

\texttt{sed} 常用于按行处理文本。

\begin{lstlisting}[language=bash]
sed 's/old/new/' file.txt
sed 's/old/new/g' file.txt
sed -i 's/old/new/g' file.txt
\end{lstlisting}

解释：
\begin{itemize}
    \item \texttt{s/old/new/}：替换每行第一个 old；
    \item \texttt{s/old/new/g}：替换所有 old；
    \item \texttt{-i}：直接修改原文件。
\end{itemize}

示例：把配置文件中的 115200 改成 9600：

\begin{lstlisting}[language=bash]
sed -i 's/115200/9600/g' uart.conf
\end{lstlisting}

\subsection{awk 文本处理}

\texttt{awk} 适合处理列数据。

\begin{lstlisting}[language=bash]
awk '{print $1}' file.txt
awk '{print $1, $3}' file.txt
awk -F ':' '{print $1}' /etc/passwd
\end{lstlisting}

示例：查看当前系统用户名：

\begin{lstlisting}[language=bash]
awk -F ':' '{print $1}' /etc/passwd
\end{lstlisting}

示例：查看磁盘使用率：

\begin{lstlisting}[language=bash]
df -h | awk '{print $1, $5, $6}'
\end{lstlisting}

\section{输入输出重定向与管道}

\subsection{标准输入、输出、错误}

Linux 中有三个标准文件描述符：

\begin{longtable}{p{4cm}p{4cm}p{6cm}}
\toprule
名称 & 编号 & 含义 \\
\midrule
stdin & 0 & 标准输入 \\
stdout & 1 & 标准输出 \\
stderr & 2 & 标准错误 \\
\bottomrule
\end{longtable}

\subsection{重定向}

\begin{lstlisting}[language=bash]
echo "hello" > a.txt
echo "world" >> a.txt
cat < a.txt
ls not_exist 2> error.log
ls /home > out.log 2> err.log
ls /home > all.log 2>&1
\end{lstlisting}

解释：
\begin{itemize}
    \item \texttt{>}：覆盖写入；
    \item \texttt{>>}：追加写入；
    \item \texttt{2>}：重定向错误输出；
    \item \texttt{2>\&1}：把错误输出合并到标准输出。
\end{itemize}

\subsection{管道}

管道 \texttt{|} 用于把前一个命令的输出作为后一个命令的输入。

\begin{lstlisting}[language=bash]
ps aux | grep ssh
dmesg | grep tty
cat file.txt | grep "error" | wc -l
\end{lstlisting}

示例：统计日志中 error 出现次数：

\begin{lstlisting}[language=bash]
grep -i "error" system.log | wc -l
\end{lstlisting}

\section{用户、用户组与权限}

\subsection{用户相关命令}

\begin{lstlisting}[language=bash]
whoami
id
who
w
su
sudo command
passwd
\end{lstlisting}

\subsection{用户管理}

\begin{lstlisting}[language=bash]
sudo useradd testuser
sudo passwd testuser
sudo userdel testuser
sudo userdel -r testuser
sudo groupadd dev
sudo usermod -aG dev testuser
\end{lstlisting}

解释：
\begin{itemize}
    \item \texttt{useradd}：添加用户；
    \item \texttt{passwd}：设置密码；
    \item \texttt{userdel -r}：删除用户及其家目录；
    \item \texttt{usermod -aG}：把用户加入组。
\end{itemize}

\subsection{文件权限}

查看文件权限：

\begin{lstlisting}[language=bash]
ls -l
\end{lstlisting}

输出示例：

\begin{lstlisting}
-rwxr-xr-- 1 user group 1024 May 10 test.sh
\end{lstlisting}

含义：
\begin{itemize}
    \item 第 1 位：文件类型，\texttt{-} 表示普通文件，\texttt{d} 表示目录；
    \item 后 9 位：权限；
    \item 每 3 位一组：所有者、所属组、其他用户；
    \item r：读，w：写，x：执行。
\end{itemize}

权限数字表示：

\begin{longtable}{p{4cm}p{4cm}}
\toprule
权限 & 数字 \\
\midrule
r & 4 \\
w & 2 \\
x & 1 \\
rwx & 7 \\
rw- & 6 \\
r-x & 5 \\
r-- & 4 \\
\bottomrule
\end{longtable}

修改权限：

\begin{lstlisting}[language=bash]
chmod 755 test.sh
chmod +x test.sh
chmod u+x test.sh
chmod g-w test.sh
chmod o-r test.sh
\end{lstlisting}

修改所有者和组：

\begin{lstlisting}[language=bash]
sudo chown user file.txt
sudo chown user:group file.txt
sudo chgrp group file.txt
\end{lstlisting}

\subsection{权限考试重点}

\begin{lstlisting}[language=bash]
chmod 755 file
\end{lstlisting}

表示：
\begin{itemize}
    \item 所有者：7 = r + w + x；
    \item 所属组：5 = r + x；
    \item 其他用户：5 = r + x。
\end{itemize}

\begin{lstlisting}[language=bash]
chmod 644 file
\end{lstlisting}

表示：
\begin{itemize}
    \item 所有者可读写；
    \item 组用户只读；
    \item 其他用户只读。
\end{itemize}

\section{进程管理}

\subsection{查看进程}

\begin{lstlisting}[language=bash]
ps
ps aux
top
htop
pidof sshd
pgrep ssh
\end{lstlisting}

常见命令：

\begin{lstlisting}[language=bash]
ps aux | grep nginx
\end{lstlisting}

\subsection{终止进程}

\begin{lstlisting}[language=bash]
kill PID
kill -9 PID
pkill process_name
killall process_name
\end{lstlisting}

常用信号：

\begin{longtable}{p{4cm}p{4cm}p{6cm}}
\toprule
信号 & 数字 & 作用 \\
\midrule
SIGTERM & 15 & 正常终止进程 \\
SIGKILL & 9 & 强制杀死进程 \\
SIGINT & 2 & Ctrl + C 产生，中断进程 \\
SIGHUP & 1 & 挂起，常用于重新加载配置 \\
\bottomrule
\end{longtable}

示例：

\begin{lstlisting}[language=bash]
ps aux | grep my_app
kill 1234
kill -9 1234
\end{lstlisting}

\subsection{前台与后台任务}

\begin{lstlisting}[language=bash]
./program
./program &
jobs
fg %1
bg %1
Ctrl + Z
Ctrl + C
\end{lstlisting}

解释：
\begin{itemize}
    \item \texttt{\&}：后台运行；
    \item \texttt{jobs}：查看后台任务；
    \item \texttt{fg}：切回前台；
    \item \texttt{bg}：后台继续运行；
    \item \texttt{Ctrl + Z}：暂停任务；
    \item \texttt{Ctrl + C}：终止任务。
\end{itemize}

\section{软件包管理}

\subsection{Debian/Ubuntu：apt}

\begin{lstlisting}[language=bash]
sudo apt update
sudo apt upgrade
sudo apt install gcc make git
sudo apt remove package_name
sudo apt purge package_name
sudo apt search package_name
sudo apt show package_name
\end{lstlisting}

示例：安装 C 语言开发工具：

\begin{lstlisting}[language=bash]
sudo apt update
sudo apt install build-essential
\end{lstlisting}

\subsection{RedHat/CentOS/Fedora：dnf 或 yum}

\begin{lstlisting}[language=bash]
sudo dnf install gcc make git
sudo dnf remove package_name
sudo dnf search package_name
\end{lstlisting}

较老系统可能使用：

\begin{lstlisting}[language=bash]
sudo yum install gcc make git
\end{lstlisting}

\section{压缩与归档}

\subsection{tar}

\begin{lstlisting}[language=bash]
tar -cvf archive.tar dir
tar -xvf archive.tar
tar -czvf archive.tar.gz dir
tar -xzvf archive.tar.gz
tar -cjvf archive.tar.bz2 dir
tar -xjvf archive.tar.bz2
\end{lstlisting}

选项含义：
\begin{itemize}
    \item \texttt{c}：创建；
    \item \texttt{x}：解压；
    \item \texttt{v}：显示过程；
    \item \texttt{f}：指定文件；
    \item \texttt{z}：gzip；
    \item \texttt{j}：bzip2。
\end{itemize}

示例：打包工程目录：

\begin{lstlisting}[language=bash]
tar -czvf project_backup.tar.gz project/
\end{lstlisting}

\subsection{zip 和 unzip}

\begin{lstlisting}[language=bash]
zip file.zip a.txt b.txt
zip -r project.zip project/
unzip project.zip
\end{lstlisting}

\section{磁盘、文件系统与挂载}

\subsection{查看磁盘信息}

\begin{lstlisting}[language=bash]
df -h
du -sh *
lsblk
fdisk -l
mount
\end{lstlisting}

解释：
\begin{itemize}
    \item \texttt{df -h}：查看文件系统空间；
    \item \texttt{du -sh}：查看文件或目录大小；
    \item \texttt{lsblk}：查看块设备；
    \item \texttt{fdisk -l}：查看磁盘分区；
    \item \texttt{mount}：查看挂载信息。
\end{itemize}

\subsection{挂载和卸载}

\begin{lstlisting}[language=bash]
sudo mount /dev/sdb1 /mnt
sudo umount /mnt
\end{lstlisting}

嵌入式中常见挂载：

\begin{lstlisting}[language=bash]
mount -t nfs 192.168.1.100:/home/nfs /mnt
\end{lstlisting}

表示通过 NFS 挂载主机目录到开发板。

\section{网络基础命令}

\subsection{查看网络配置}

\begin{lstlisting}[language=bash]
ip addr
ip link
ifconfig
hostname
hostname -I
\end{lstlisting}

推荐新命令：

\begin{lstlisting}[language=bash]
ip addr show
ip route show
\end{lstlisting}

\subsection{测试网络}

\begin{lstlisting}[language=bash]
ping 8.8.8.8
ping www.baidu.com
traceroute www.baidu.com
curl www.example.com
wget http://example.com/file.tar.gz
\end{lstlisting}

\subsection{端口和连接}

\begin{lstlisting}[language=bash]
netstat -tulnp
ss -tulnp
lsof -i :8080
\end{lstlisting}

示例：查看 8080 端口被谁占用：

\begin{lstlisting}[language=bash]
sudo lsof -i :8080
\end{lstlisting}

\subsection{SSH 远程登录}

\begin{lstlisting}[language=bash]
ssh user@192.168.1.10
scp file.txt user@192.168.1.10:/home/user/
scp user@192.168.1.10:/home/user/a.txt .
\end{lstlisting}

嵌入式开发中常用 SSH 登录开发板：

\begin{lstlisting}[language=bash]
ssh root@192.168.1.100
\end{lstlisting}

\section{环境变量}

\subsection{查看环境变量}

\begin{lstlisting}[language=bash]
env
printenv
echo $PATH
echo $HOME
echo $USER
\end{lstlisting}

\subsection{设置环境变量}

临时设置：

\begin{lstlisting}[language=bash]
export MYVAR=hello
echo $MYVAR
\end{lstlisting}

永久设置可写入：

\begin{lstlisting}[language=bash]
~/.bashrc
~/.profile
/etc/profile
\end{lstlisting}

示例：

\begin{lstlisting}[language=bash]
echo 'export PATH=$PATH:/opt/toolchain/bin' >> ~/.bashrc
source ~/.bashrc
\end{lstlisting}

嵌入式交叉编译经常设置工具链路径：

\begin{lstlisting}[language=bash]
export PATH=$PATH:/opt/arm-linux-gcc/bin
export CROSS_COMPILE=arm-linux-gnueabihf-
export ARCH=arm
\end{lstlisting}

\section{Vim 基础}

\subsection{Vim 三种常用模式}

\begin{itemize}
    \item 普通模式：移动光标、复制、删除；
    \item 插入模式：输入文本；
    \item 命令模式：保存、退出、查找。
\end{itemize}

\subsection{常用操作}

\begin{longtable}{p{4cm}p{10cm}}
\toprule
命令 & 作用 \\
\midrule
i & 进入插入模式 \\
Esc & 返回普通模式 \\
:w & 保存 \\
:q & 退出 \\
:wq & 保存并退出 \\
:q! & 强制退出不保存 \\
dd & 删除一行 \\
yy & 复制一行 \\
p & 粘贴 \\
/word & 查找 word \\
n & 下一个匹配 \\
u & 撤销 \\
\bottomrule
\end{longtable}

\section{Shell 脚本基础}

\subsection{第一个 Shell 脚本}

创建 \texttt{hello.sh}：

\begin{lstlisting}[language=bash]
#!/bin/bash

echo "Hello Linux"
\end{lstlisting}

执行：

\begin{lstlisting}[language=bash]
chmod +x hello.sh
./hello.sh
\end{lstlisting}

解释：
\begin{itemize}
    \item \texttt{\#!/bin/bash}：指定脚本解释器；
    \item \texttt{chmod +x}：添加可执行权限。
\end{itemize}

\subsection{变量}

\begin{lstlisting}[language=bash]
#!/bin/bash

name="Linux"
echo "Hello $name"
echo "Hello ${name}"
\end{lstlisting}

注意：
\begin{itemize}
    \item 变量赋值时等号两边不能有空格；
    \item 使用变量时加 \texttt{\$}。
\end{itemize}

\subsection{命令替换}

\begin{lstlisting}[language=bash]
#!/bin/bash

now=$(date)
echo "Current time: $now"

files=$(ls)
echo "$files"
\end{lstlisting}

\subsection{位置参数}

\begin{lstlisting}[language=bash]
#!/bin/bash

echo "Script name: $0"
echo "First arg: $1"
echo "Second arg: $2"
echo "All args: $@"
echo "Arg count: $#"
\end{lstlisting}

执行：

\begin{lstlisting}[language=bash]
./test.sh hello world
\end{lstlisting}

输出中：
\begin{itemize}
    \item \texttt{\$0}：脚本名；
    \item \texttt{\$1}：第一个参数；
    \item \texttt{\$2}：第二个参数；
    \item \texttt{\$@}：所有参数；
    \item \texttt{\$\#}：参数个数。
\end{itemize}

\subsection{条件判断}

\begin{lstlisting}[language=bash]
#!/bin/bash

num=10

if [ $num -gt 5 ]; then
    echo "num > 5"
else
    echo "num <= 5"
fi
\end{lstlisting}

常用数字比较：

\begin{longtable}{p{4cm}p{10cm}}
\toprule
表达式 & 含义 \\
\midrule
-eq & 等于 \\
-ne & 不等于 \\
-gt & 大于 \\
-lt & 小于 \\
-ge & 大于等于 \\
-le & 小于等于 \\
\bottomrule
\end{longtable}

字符串比较：

\begin{lstlisting}[language=bash]
#!/bin/bash

str="yes"

if [ "$str" = "yes" ]; then
    echo "OK"
fi
\end{lstlisting}

文件判断：

\begin{longtable}{p{4cm}p{10cm}}
\toprule
表达式 & 含义 \\
\midrule
-e file & 文件存在 \\
-f file & 普通文件 \\
-d dir & 目录 \\
-r file & 可读 \\
-w file & 可写 \\
-x file & 可执行 \\
\bottomrule
\end{longtable}

示例：

\begin{lstlisting}[language=bash]
#!/bin/bash

if [ -f "main.c" ]; then
    echo "main.c exists"
else
    echo "main.c not found"
fi
\end{lstlisting}

\subsection{case 分支}

\begin{lstlisting}[language=bash]
#!/bin/bash

case "$1" in
    start)
        echo "start service"
        ;;
    stop)
        echo "stop service"
        ;;
    restart)
        echo "restart service"
        ;;
    *)
        echo "Usage: $0 {start|stop|restart}"
        ;;
esac
\end{lstlisting}

\subsection{for 循环}

\begin{lstlisting}[language=bash]
#!/bin/bash

for file in *.c
do
    echo "C file: $file"
done
\end{lstlisting}

数字循环：

\begin{lstlisting}[language=bash]
#!/bin/bash

for i in {1..5}
do
    echo $i
done
\end{lstlisting}

\subsection{while 循环}

\begin{lstlisting}[language=bash]
#!/bin/bash

count=1

while [ $count -le 5 ]
do
    echo $count
    count=$((count + 1))
done
\end{lstlisting}

\subsection{函数}

\begin{lstlisting}[language=bash]
#!/bin/bash

say_hello() {
    echo "Hello $1"
}

say_hello Linux
say_hello ROS2
\end{lstlisting}

\subsection{退出状态码}

每条命令执行后都有返回值，\texttt{\$?} 表示上一条命令状态：

\begin{lstlisting}[language=bash]
ls /home
echo $?

ls /not_exist
echo $?
\end{lstlisting}

一般：
\begin{itemize}
    \item 0 表示成功；
    \item 非 0 表示失败。
\end{itemize}

\subsection{实用脚本示例：自动备份目录}

\begin{lstlisting}[language=bash]
#!/bin/bash

src_dir="$1"
backup_dir="$2"

if [ $# -ne 2 ]; then
    echo "Usage: $0 src_dir backup_dir"
    exit 1
fi

if [ ! -d "$src_dir" ]; then
    echo "Source directory not found"
    exit 2
fi

mkdir -p "$backup_dir"

time_str=$(date +%Y%m%d_%H%M%S)
backup_file="$backup_dir/backup_$time_str.tar.gz"

tar -czvf "$backup_file" "$src_dir"

echo "Backup saved to $backup_file"
\end{lstlisting}

执行：

\begin{lstlisting}[language=bash]
chmod +x backup.sh
./backup.sh project backups
\end{lstlisting}

\subsection{实用脚本示例：检查程序是否运行}

\begin{lstlisting}[language=bash]
#!/bin/bash

process_name="$1"

if [ $# -ne 1 ]; then
    echo "Usage: $0 process_name"
    exit 1
fi

pid=$(pgrep "$process_name")

if [ -n "$pid" ]; then
    echo "$process_name is running, pid=$pid"
else
    echo "$process_name is not running"
fi
\end{lstlisting}

\section{编译、Makefile 与调试}

\subsection{gcc 编译 C 程序}

假设有 \texttt{hello.c}：

\begin{lstlisting}[language=C]
#include <stdio.h>

int main(void)
{
    printf("Hello Linux\n");
    return 0;
}
\end{lstlisting}

编译：

\begin{lstlisting}[language=bash]
gcc hello.c -o hello
./hello
\end{lstlisting}

带警告和调试信息：

\begin{lstlisting}[language=bash]
gcc -Wall -g hello.c -o hello
\end{lstlisting}

常用选项：
\begin{itemize}
    \item \texttt{-o}：指定输出文件；
    \item \texttt{-Wall}：打开常见警告；
    \item \texttt{-g}：生成调试信息；
    \item \texttt{-I}：指定头文件目录；
    \item \texttt{-L}：指定库目录；
    \item \texttt{-l}：链接库。
\end{itemize}

示例：

\begin{lstlisting}[language=bash]
gcc main.c -I./include -L./lib -lmylib -o app
\end{lstlisting}

\subsection{Makefile 基础}

简单 Makefile：

\begin{lstlisting}[language=make]
app: main.o add.o
	gcc main.o add.o -o app

main.o: main.c
	gcc -c main.c

add.o: add.c
	gcc -c add.c

clean:
	rm -f *.o app
\end{lstlisting}

注意：命令前面必须是 Tab，不是空格。

执行：

\begin{lstlisting}[language=bash]
make
make clean
\end{lstlisting}

更常见写法：

\begin{lstlisting}[language=make]
CC = gcc
CFLAGS = -Wall -g
TARGET = app
OBJS = main.o add.o

$(TARGET): $(OBJS)
	$(CC) $(OBJS) -o $(TARGET)

%.o: %.c
	$(CC) $(CFLAGS) -c $< -o $@

clean:
	rm -f $(OBJS) $(TARGET)
\end{lstlisting}

\subsection{gdb 调试}

编译时添加 \texttt{-g}：

\begin{lstlisting}[language=bash]
gcc -g main.c -o app
gdb ./app
\end{lstlisting}

gdb 常用命令：

\begin{longtable}{p{4cm}p{10cm}}
\toprule
命令 & 作用 \\
\midrule
break main & 在 main 函数设置断点 \\
b 10 & 在第 10 行设置断点 \\
run & 运行程序 \\
next & 单步执行，不进入函数 \\
step & 单步执行，进入函数 \\
continue & 继续运行 \\
print var & 打印变量 \\
backtrace & 查看函数调用栈 \\
quit & 退出 gdb \\
\bottomrule
\end{longtable}

示例：

\begin{lstlisting}
(gdb) break main
(gdb) run
(gdb) next
(gdb) print i
(gdb) continue
(gdb) quit
\end{lstlisting}

\section{Linux 系统管理}

\subsection{systemd 服务管理}

\begin{lstlisting}[language=bash]
systemctl status ssh
sudo systemctl start ssh
sudo systemctl stop ssh
sudo systemctl restart ssh
sudo systemctl enable ssh
sudo systemctl disable ssh
\end{lstlisting}

查看日志：

\begin{lstlisting}[language=bash]
journalctl
journalctl -u ssh
journalctl -f
journalctl -u ssh -f
\end{lstlisting}

\subsection{定时任务 cron}

编辑当前用户定时任务：

\begin{lstlisting}[language=bash]
crontab -e
\end{lstlisting}

查看：

\begin{lstlisting}[language=bash]
crontab -l
\end{lstlisting}

格式：

\begin{lstlisting}
分钟 小时 日期 月份 星期 命令
\end{lstlisting}

例子：每天凌晨 2 点执行备份：

\begin{lstlisting}
0 2 * * * /home/user/backup.sh
\end{lstlisting}

每 5 分钟执行一次：

\begin{lstlisting}
*/5 * * * * /home/user/check.sh
\end{lstlisting}

\section{嵌入式 Linux 基础}

\subsection{嵌入式 Linux 启动流程}

典型启动流程：

\begin{enumerate}
    \item 上电复位；
    \item BootROM 执行；
    \item 加载 Bootloader，例如 U-Boot；
    \item Bootloader 初始化 DDR、Flash、串口等硬件；
    \item Bootloader 加载 Linux Kernel；
    \item Kernel 解压并启动；
    \item Kernel 挂载根文件系统 rootfs；
    \item 启动第一个用户空间进程 init 或 systemd；
    \item 启动各种服务和应用程序。
\end{enumerate}

简化流程：

\begin{lstlisting}
BootROM -> U-Boot -> Linux Kernel -> RootFS -> init/systemd -> App
\end{lstlisting}

\subsection{U-Boot 常用命令}

\begin{lstlisting}
printenv
setenv bootargs console=ttyS0,115200 root=/dev/mmcblk0p2 rw
saveenv
boot
bootm
tftpboot
mmc list
mmc dev 0
fatls mmc 0:1
fatload mmc 0:1 0x80000000 zImage
\end{lstlisting}

常见环境变量：
\begin{itemize}
    \item \texttt{bootargs}：传给 Linux 内核的启动参数；
    \item \texttt{bootcmd}：自动启动命令；
    \item \texttt{ethaddr}：网卡 MAC 地址；
    \item \texttt{ipaddr}：开发板 IP；
    \item \texttt{serverip}：服务器 IP。
\end{itemize}

\subsection{内核启动参数}

示例：

\begin{lstlisting}
console=ttyS0,115200 root=/dev/mmcblk0p2 rw rootwait
\end{lstlisting}

含义：
\begin{itemize}
    \item \texttt{console=ttyS0,115200}：串口控制台；
    \item \texttt{root=/dev/mmcblk0p2}：根文件系统位置；
    \item \texttt{rw}：以读写方式挂载；
    \item \texttt{rootwait}：等待根设备准备好。
\end{itemize}

\subsection{设备文件}

Linux 中设备也表现为文件，通常在 \texttt{/dev} 目录下。

常见设备：
\begin{longtable}{p{4cm}p{10cm}}
\toprule
设备文件 & 含义 \\
\midrule
/dev/ttyS0 & 串口设备 \\
/dev/ttyUSB0 & USB 转串口设备 \\
/dev/mmcblk0 & eMMC 或 SD 卡 \\
/dev/sda & SATA 或 USB 存储设备 \\
/dev/i2c-0 & I2C 总线设备 \\
/dev/spidev0.0 & SPI 设备 \\
/dev/fb0 & Framebuffer 显示设备 \\
/dev/input/event0 & 输入设备 \\
\bottomrule
\end{longtable}

查看串口设备：

\begin{lstlisting}[language=bash]
ls /dev/ttyS*
ls /dev/ttyUSB*
dmesg | grep tty
\end{lstlisting}

\subsection{proc 文件系统}

\texttt{/proc} 是虚拟文件系统，用于查看内核和进程信息。

\begin{lstlisting}[language=bash]
cat /proc/cpuinfo
cat /proc/meminfo
cat /proc/version
cat /proc/cmdline
cat /proc/interrupts
cat /proc/filesystems
\end{lstlisting}

嵌入式常用：

\begin{lstlisting}[language=bash]
cat /proc/cmdline
cat /proc/interrupts
cat /proc/meminfo
\end{lstlisting}

\subsection{sys 文件系统}

\texttt{/sys} 用于表示设备、驱动、总线等内核对象。

常见路径：

\begin{lstlisting}[language=bash]
/sys/class
/sys/bus
/sys/devices
/sys/class/gpio
/sys/class/leds
/sys/class/pwm
\end{lstlisting}

GPIO 旧式 sysfs 操作示例：

\begin{lstlisting}[language=bash]
echo 17 > /sys/class/gpio/export
echo out > /sys/class/gpio/gpio17/direction
echo 1 > /sys/class/gpio/gpio17/value
echo 0 > /sys/class/gpio/gpio17/value
echo 17 > /sys/class/gpio/unexport
\end{lstlisting}

注意：较新的 Linux 内核更推荐使用 \texttt{libgpiod}。

\subsection{设备树 Device Tree}

设备树用于描述硬件信息，让内核知道板级硬件资源。

常见文件：
\begin{itemize}
    \item \texttt{.dts}：设备树源文件；
    \item \texttt{.dtsi}：设备树包含文件；
    \item \texttt{.dtb}：编译后的设备树二进制文件。
\end{itemize}

设备树示例：

\begin{lstlisting}
uart0: serial@101f1000 {
    compatible = "arm,pl011";
    reg = <0x101f1000 0x1000>;
    interrupts = <5>;
    status = "okay";
};
\end{lstlisting}

常见属性：
\begin{itemize}
    \item \texttt{compatible}：用于驱动匹配；
    \item \texttt{reg}：寄存器地址和大小；
    \item \texttt{interrupts}：中断号；
    \item \texttt{status}：设备是否启用。
\end{itemize}

\subsection{内核模块}

查看模块：

\begin{lstlisting}[language=bash]
lsmod
\end{lstlisting}

加载模块：

\begin{lstlisting}[language=bash]
sudo insmod hello.ko
sudo modprobe driver_name
\end{lstlisting}

卸载模块：

\begin{lstlisting}[language=bash]
sudo rmmod hello
\end{lstlisting}

查看内核日志：

\begin{lstlisting}[language=bash]
dmesg
dmesg | tail
dmesg -w
\end{lstlisting}

\subsection{交叉编译}

交叉编译是指在一种平台上编译另一种平台运行的程序。

例如在 x86 PC 上编译 ARM 程序：

\begin{lstlisting}[language=bash]
arm-linux-gnueabihf-gcc main.c -o app
\end{lstlisting}

常见环境变量：

\begin{lstlisting}[language=bash]
export ARCH=arm
export CROSS_COMPILE=arm-linux-gnueabihf-
\end{lstlisting}

编译内核示例：

\begin{lstlisting}[language=bash]
make ARCH=arm CROSS_COMPILE=arm-linux-gnueabihf- menuconfig
make ARCH=arm CROSS_COMPILE=arm-linux-gnueabihf- zImage
make ARCH=arm CROSS_COMPILE=arm-linux-gnueabihf- dtbs
\end{lstlisting}

\subsection{NFS 根文件系统}

开发阶段常用 NFS 挂载根文件系统，方便调试。

主机端安装 NFS：

\begin{lstlisting}[language=bash]
sudo apt install nfs-kernel-server
\end{lstlisting}

配置 \texttt{/etc/exports}：

\begin{lstlisting}
/home/user/nfsroot *(rw,sync,no_root_squash,no_subtree_check)
\end{lstlisting}

重启服务：

\begin{lstlisting}[language=bash]
sudo exportfs -ra
sudo systemctl restart nfs-kernel-server
\end{lstlisting}

开发板挂载：

\begin{lstlisting}[language=bash]
mount -t nfs 192.168.1.10:/home/user/nfsroot /mnt
\end{lstlisting}

\section{Linux C 编程基础}

\subsection{文件 I/O}

系统调用方式：

\begin{lstlisting}[language=C]
#include <stdio.h>
#include <fcntl.h>
#include <unistd.h>

int main(void)
{
    int fd;
    char buf[100];

    fd = open("test.txt", O_RDONLY);
    if (fd < 0) {
        perror("open");
        return 1;
    }

    int n = read(fd, buf, sizeof(buf) - 1);
    if (n < 0) {
        perror("read");
        close(fd);
        return 1;
    }

    buf[n] = '\0';
    printf("%s\n", buf);

    close(fd);
    return 0;
}
\end{lstlisting}

常见系统调用：
\begin{itemize}
    \item \texttt{open}：打开文件；
    \item \texttt{read}：读取文件；
    \item \texttt{write}：写文件；
    \item \texttt{close}：关闭文件；
    \item \texttt{ioctl}：设备控制，驱动开发常见。
\end{itemize}

\subsection{进程创建 fork}

\begin{lstlisting}[language=C]
#include <stdio.h>
#include <unistd.h>

int main(void)
{
    pid_t pid = fork();

    if (pid < 0) {
        perror("fork");
        return 1;
    } else if (pid == 0) {
        printf("Child process\n");
    } else {
        printf("Parent process, child pid=%d\n", pid);
    }

    return 0;
}
\end{lstlisting}

\subsection{线程 pthread}

\begin{lstlisting}[language=C]
#include <stdio.h>
#include <pthread.h>

void *thread_func(void *arg)
{
    printf("Hello from thread\n");
    return NULL;
}

int main(void)
{
    pthread_t tid;

    pthread_create(&tid, NULL, thread_func, NULL);
    pthread_join(tid, NULL);

    return 0;
}
\end{lstlisting}

编译：

\begin{lstlisting}[language=bash]
gcc thread.c -o thread -lpthread
\end{lstlisting}

\section{ROS 2 基础知识}

\subsection{ROS 2 是什么}

ROS 2 全称 Robot Operating System 2，不是传统意义上的操作系统，而是机器人软件开发框架。

ROS 2 提供：
\begin{itemize}
    \item 节点通信；
    \item 话题发布订阅；
    \item 服务请求响应；
    \item 动作通信；
    \item 参数管理；
    \item 坐标变换；
    \item 日志系统；
    \item 包管理和构建工具。
\end{itemize}

ROS 2 与 ROS 1 的重要区别：
\begin{itemize}
    \item ROS 2 基于 DDS 通信；
    \item 支持实时性更好；
    \item 支持多机器人通信；
    \item 支持 QoS；
    \item 支持生命周期节点；
    \item 不再依赖 ROS Master。
\end{itemize}

\subsection{ROS 2 核心概念}

\begin{longtable}{p{4cm}p{10cm}}
\toprule
概念 & 含义 \\
\midrule
Node & 节点，一个独立功能模块 \\
Topic & 话题，用于发布订阅通信 \\
Publisher & 发布者，向 Topic 发送消息 \\
Subscriber & 订阅者，从 Topic 接收消息 \\
Service & 服务，请求-响应模式 \\
Client & 服务客户端，发送请求 \\
Server & 服务端，处理请求 \\
Action & 动作，适合长时间任务 \\
Message & 消息类型 \\
Parameter & 参数，节点运行时配置 \\
Launch & 启动文件，用于同时启动多个节点 \\
Workspace & 工作空间，存放 ROS 2 工程 \\
Package & 功能包，ROS 2 基本组织单位 \\
QoS & 服务质量策略 \\
DDS & ROS 2 底层通信机制 \\
\bottomrule
\end{longtable}

\subsection{ROS 2 工作空间}

典型工作空间结构：

\begin{lstlisting}
ros2_ws/
├── src/
│   └── my_package/
├── build/
├── install/
└── log/
\end{lstlisting}

创建工作空间：

\begin{lstlisting}[language=bash]
mkdir -p ~/ros2_ws/src
cd ~/ros2_ws
colcon build
source install/setup.bash
\end{lstlisting}

每次打开新终端后需要：

\begin{lstlisting}[language=bash]
source /opt/ros/humble/setup.bash
source ~/ros2_ws/install/setup.bash
\end{lstlisting}

也可以写入 \texttt{~/.bashrc}：

\begin{lstlisting}[language=bash]
echo "source /opt/ros/humble/setup.bash" >> ~/.bashrc
echo "source ~/ros2_ws/install/setup.bash" >> ~/.bashrc
source ~/.bashrc
\end{lstlisting}

\subsection{ROS 2 常用命令}

\subsubsection{查看节点}

\begin{lstlisting}[language=bash]
ros2 node list
ros2 node info /node_name
\end{lstlisting}

\subsubsection{查看话题}

\begin{lstlisting}[language=bash]
ros2 topic list
ros2 topic info /topic_name
ros2 topic echo /topic_name
ros2 topic hz /topic_name
ros2 topic pub /chatter std_msgs/msg/String "{data: 'hello'}"
\end{lstlisting}

示例：

\begin{lstlisting}[language=bash]
ros2 topic pub /cmd_vel geometry_msgs/msg/Twist \
"{linear: {x: 0.2, y: 0.0, z: 0.0}, angular: {x: 0.0, y: 0.0, z: 0.5}}"
\end{lstlisting}

\subsubsection{查看服务}

\begin{lstlisting}[language=bash]
ros2 service list
ros2 service type /service_name
ros2 service call /service_name service_type "{request}"
\end{lstlisting}

示例：

\begin{lstlisting}[language=bash]
ros2 service call /add_two_ints example_interfaces/srv/AddTwoInts "{a: 1, b: 2}"
\end{lstlisting}

\subsubsection{查看参数}

\begin{lstlisting}[language=bash]
ros2 param list
ros2 param get /node_name parameter_name
ros2 param set /node_name parameter_name value
ros2 param dump /node_name
\end{lstlisting}

\subsubsection{查看接口类型}

\begin{lstlisting}[language=bash]
ros2 interface list
ros2 interface show std_msgs/msg/String
ros2 interface show geometry_msgs/msg/Twist
\end{lstlisting}

\subsubsection{运行节点}

\begin{lstlisting}[language=bash]
ros2 run package_name executable_name
\end{lstlisting}

示例：

\begin{lstlisting}[language=bash]
ros2 run turtlesim turtlesim_node
ros2 run turtlesim turtle_teleop_key
\end{lstlisting}

\subsection{创建 ROS 2 Python 功能包}

创建包：

\begin{lstlisting}[language=bash]
cd ~/ros2_ws/src
ros2 pkg create my_py_pkg --build-type ament_python --dependencies rclpy std_msgs
\end{lstlisting}

目录结构大致为：

\begin{lstlisting}
my_py_pkg/
├── my_py_pkg/
│   ├── __init__.py
├── package.xml
├── setup.py
└── setup.cfg
\end{lstlisting}

\subsection{Python 发布者示例}

文件：\texttt{my\_py\_pkg/talker.py}

\begin{lstlisting}[language=Python]
import rclpy
from rclpy.node import Node
from std_msgs.msg import String

class Talker(Node):
    def __init__(self):
        super().__init__('talker')
        self.publisher = self.create_publisher(String, 'chatter', 10)
        self.timer = self.create_timer(1.0, self.timer_callback)
        self.count = 0

    def timer_callback(self):
        msg = String()
        msg.data = f'Hello ROS 2: {self.count}'
        self.publisher.publish(msg)
        self.get_logger().info(f'Publish: {msg.data}')
        self.count += 1

def main(args=None):
    rclpy.init(args=args)
    node = Talker()
    rclpy.spin(node)
    node.destroy_node()
    rclpy.shutdown()

if __name__ == '__main__':
    main()
\end{lstlisting}

在 \texttt{setup.py} 中添加入口：

\begin{lstlisting}[language=Python]
entry_points={
    'console_scripts': [
        'talker = my_py_pkg.talker:main',
    ],
},
\end{lstlisting}

编译运行：

\begin{lstlisting}[language=bash]
cd ~/ros2_ws
colcon build
source install/setup.bash
ros2 run my_py_pkg talker
\end{lstlisting}

\subsection{Python 订阅者示例}

文件：\texttt{my\_py\_pkg/listener.py}

\begin{lstlisting}[language=Python]
import rclpy
from rclpy.node import Node
from std_msgs.msg import String

class Listener(Node):
    def __init__(self):
        super().__init__('listener')
        self.subscription = self.create_subscription(
            String,
            'chatter',
            self.callback,
            10
        )

    def callback(self, msg):
        self.get_logger().info(f'I heard: {msg.data}')

def main(args=None):
    rclpy.init(args=args)
    node = Listener()
    rclpy.spin(node)
    node.destroy_node()
    rclpy.shutdown()

if __name__ == '__main__':
    main()
\end{lstlisting}

在 \texttt{setup.py} 中添加：

\begin{lstlisting}[language=Python]
entry_points={
    'console_scripts': [
        'talker = my_py_pkg.talker:main',
        'listener = my_py_pkg.listener:main',
    ],
},
\end{lstlisting}

编译运行：

\begin{lstlisting}[language=bash]
cd ~/ros2_ws
colcon build
source install/setup.bash

ros2 run my_py_pkg talker
ros2 run my_py_pkg listener
\end{lstlisting}

\subsection{创建 ROS 2 C++ 功能包}

\begin{lstlisting}[language=bash]
cd ~/ros2_ws/src
ros2 pkg create my_cpp_pkg --build-type ament_cmake --dependencies rclcpp std_msgs
\end{lstlisting}

\subsection{C++ 发布者示例}

文件：\texttt{src/talker.cpp}

\begin{lstlisting}[language=C++]
#include <chrono>
#include <memory>
#include <string>

#include "rclcpp/rclcpp.hpp"
#include "std_msgs/msg/string.hpp"

using namespace std::chrono_literals;

class Talker : public rclcpp::Node
{
public:
    Talker() : Node("cpp_talker"), count_(0)
    {
        publisher_ = this->create_publisher<std_msgs::msg::String>("chatter", 10);

        timer_ = this->create_wall_timer(
            1s,
            std::bind(&Talker::timer_callback, this)
        );
    }

private:
    void timer_callback()
    {
        auto msg = std_msgs::msg::String();
        msg.data = "Hello ROS 2 C++: " + std::to_string(count_++);
        publisher_->publish(msg);
        RCLCPP_INFO(this->get_logger(), "Publish: '%s'", msg.data.c_str());
    }

    rclcpp::Publisher<std_msgs::msg::String>::SharedPtr publisher_;
    rclcpp::TimerBase::SharedPtr timer_;
    size_t count_;
};

int main(int argc, char * argv[])
{
    rclcpp::init(argc, argv);
    rclcpp::spin(std::make_shared<Talker>());
    rclcpp::shutdown();
    return 0;
}
\end{lstlisting}

\subsection{CMakeLists.txt 示例}

\begin{lstlisting}[language=Cmake]
cmake_minimum_required(VERSION 3.8)
project(my_cpp_pkg)

find_package(ament_cmake REQUIRED)
find_package(rclcpp REQUIRED)
find_package(std_msgs REQUIRED)

add_executable(talker src/talker.cpp)
ament_target_dependencies(talker rclcpp std_msgs)

install(TARGETS
  talker
  DESTINATION lib/${PROJECT_NAME}
)

ament_package()
\end{lstlisting}

编译运行：

\begin{lstlisting}[language=bash]
cd ~/ros2_ws
colcon build
source install/setup.bash
ros2 run my_cpp_pkg talker
\end{lstlisting}

\section{ROS 2 通信机制}

\subsection{Topic：发布订阅}

Topic 适合连续数据流，例如：
\begin{itemize}
    \item 摄像头图像；
    \item 激光雷达数据；
    \item 速度控制命令；
    \item 机器人位姿；
    \item IMU 数据。
\end{itemize}

特点：
\begin{itemize}
    \item 异步通信；
    \item 多个发布者和多个订阅者；
    \item 不要求发送方等待接收方回复。
\end{itemize}

\subsection{Service：请求响应}

Service 适合短时间请求，例如：
\begin{itemize}
    \item 查询状态；
    \item 开关设备；
    \item 重置机器人；
    \item 请求一次计算。
\end{itemize}

特点：
\begin{itemize}
    \item 同步或异步请求响应；
    \item 客户端发送请求；
    \item 服务端返回结果。
\end{itemize}

\subsection{Action：长时间任务}

Action 适合耗时任务，例如：
\begin{itemize}
    \item 导航到目标点；
    \item 机械臂执行轨迹；
    \item 自动充电；
    \item 巡航任务。
\end{itemize}

Action 包含：
\begin{itemize}
    \item Goal：目标；
    \item Feedback：执行过程反馈；
    \item Result：最终结果。
\end{itemize}

\subsection{Parameter：参数}

参数用于配置节点。

命令示例：

\begin{lstlisting}[language=bash]
ros2 param list
ros2 param get /my_node speed
ros2 param set /my_node speed 1.0
\end{lstlisting}

Python 节点中声明参数：

\begin{lstlisting}[language=Python]
self.declare_parameter('speed', 1.0)
speed = self.get_parameter('speed').value
\end{lstlisting}

\subsection{QoS 服务质量}

ROS 2 QoS 用于控制通信质量。

常见策略：
\begin{itemize}
    \item Reliability：可靠性；
    \item Durability：持久性；
    \item History：历史缓存；
    \item Depth：队列深度；
    \item Deadline：截止时间；
    \item Lifespan：消息寿命。
\end{itemize}

常见可靠性：
\begin{itemize}
    \item Reliable：可靠传输，适合重要数据；
    \item Best Effort：尽力传输，适合传感器高速数据。
\end{itemize}

示例：
\begin{itemize}
    \item 激光雷达、摄像头：常用 Best Effort；
    \item 控制命令、状态：常用 Reliable。
\end{itemize}

\section{ROS 2 Launch 文件}

Launch 文件用于一次启动多个节点。

Python launch 示例：

\begin{lstlisting}[language=Python]
from launch import LaunchDescription
from launch_ros.actions import Node

def generate_launch_description():
    return LaunchDescription([
        Node(
            package='my_py_pkg',
            executable='talker',
            name='talker_node'
        ),
        Node(
            package='my_py_pkg',
            executable='listener',
            name='listener_node'
        )
    ])
\end{lstlisting}

运行：

\begin{lstlisting}[language=bash]
ros2 launch my_py_pkg demo.launch.py
\end{lstlisting}

安装 launch 文件时，在 \texttt{setup.py} 中通常需要配置数据文件。

\section{ROS 2 TF2 坐标变换基础}

TF2 用于维护机器人各坐标系之间的关系。

常见坐标系：
\begin{itemize}
    \item \texttt{map}：地图坐标系；
    \item \texttt{odom}：里程计坐标系；
    \item \texttt{base\_link}：机器人本体坐标系；
    \item \texttt{laser}：雷达坐标系；
    \item \texttt{camera\_link}：相机坐标系。
\end{itemize}

常见关系：

\begin{lstlisting}
map -> odom -> base_link -> laser
                  |
                  -> camera_link
\end{lstlisting}

查看 TF：

\begin{lstlisting}[language=bash]
ros2 run tf2_tools view_frames
ros2 run tf2_ros tf2_echo map base_link
\end{lstlisting}

发布静态 TF：

\begin{lstlisting}[language=bash]
ros2 run tf2_ros static_transform_publisher \
0 0 0 0 0 0 base_link laser
\end{lstlisting}

\section{ROS 2 Bag 录包与回放}

录制所有话题：

\begin{lstlisting}[language=bash]
ros2 bag record -a
\end{lstlisting}

录制指定话题：

\begin{lstlisting}[language=bash]
ros2 bag record /chatter /cmd_vel
\end{lstlisting}

查看 bag 信息：

\begin{lstlisting}[language=bash]
ros2 bag info bag_directory
\end{lstlisting}

回放：

\begin{lstlisting}[language=bash]
ros2 bag play bag_directory
\end{lstlisting}

常用于：
\begin{itemize}
    \item 记录传感器数据；
    \item 离线调试算法；
    \item 比赛和实验数据复现。
\end{itemize}

\section{ROS 2 常见排错}

\subsection{找不到包}

现象：

\begin{lstlisting}
Package 'xxx' not found
\end{lstlisting}

解决：

\begin{lstlisting}[language=bash]
cd ~/ros2_ws
colcon build
source install/setup.bash
ros2 pkg list | grep xxx
\end{lstlisting}

\subsection{节点运行不了}

检查：
\begin{lstlisting}[language=bash]
ros2 pkg executables package_name
ros2 run package_name executable_name
\end{lstlisting}

Python 包还要检查 \texttt{setup.py} 中的 \texttt{entry\_points}。

\subsection{话题没有数据}

检查：
\begin{lstlisting}[language=bash]
ros2 topic list
ros2 topic info /topic_name
ros2 topic echo /topic_name
ros2 node list
ros2 node info /node_name
\end{lstlisting}

可能原因：
\begin{itemize}
    \item 节点没启动；
    \item 话题名不一致；
    \item 消息类型不一致；
    \item QoS 不匹配；
    \item 没有 source 环境。
\end{itemize}

\subsection{多机通信失败}

检查：
\begin{lstlisting}[language=bash]
echo $ROS_DOMAIN_ID
echo $RMW_IMPLEMENTATION
ip addr
ping other_machine_ip
\end{lstlisting}

两台机器需要：
\begin{itemize}
    \item 网络互通；
    \item ROS 2 环境一致或兼容；
    \item \texttt{ROS\_DOMAIN\_ID} 一致；
    \item 防火墙不阻断 DDS 通信。
\end{itemize}

\section{Linux 与 ROS 2 期末速记}

\subsection{Linux 命令速记}

\begin{longtable}{p{4cm}p{10cm}}
\toprule
任务 & 命令 \\
\midrule
查看当前目录 & pwd \\
进入目录 & cd dir \\
查看文件 & ls -l \\
创建目录 & mkdir dir \\
复制文件 & cp a b \\
复制目录 & cp -r dir1 dir2 \\
移动或重命名 & mv a b \\
删除文件 & rm file \\
删除目录 & rm -r dir \\
查看文件内容 & cat file \\
分页查看 & less file \\
查看结尾 & tail -n 20 file \\
实时查看日志 & tail -f log \\
查找文本 & grep "word" file \\
查找文件 & find . -name "*.c" \\
查看进程 & ps aux \\
杀死进程 & kill PID \\
查看磁盘 & df -h \\
查看目录大小 & du -sh dir \\
查看网络 & ip addr \\
远程登录 & ssh user@ip \\
复制远程文件 & scp file user@ip:/path \\
打包压缩 & tar -czvf a.tar.gz dir \\
解压 & tar -xzvf a.tar.gz \\
\bottomrule
\end{longtable}

\subsection{ROS 2 命令速记}

\begin{longtable}{p{4cm}p{10cm}}
\toprule
任务 & 命令 \\
\midrule
查看节点 & ros2 node list \\
查看节点信息 & ros2 node info /node \\
查看话题 & ros2 topic list \\
查看话题数据 & ros2 topic echo /topic \\
查看话题频率 & ros2 topic hz /topic \\
发布话题 & ros2 topic pub /topic type data \\
查看服务 & ros2 service list \\
调用服务 & ros2 service call /srv type data \\
查看参数 & ros2 param list \\
获取参数 & ros2 param get /node param \\
设置参数 & ros2 param set /node param value \\
查看接口 & ros2 interface list \\
运行节点 & ros2 run pkg exe \\
启动 launch & ros2 launch pkg file.launch.py \\
创建包 & ros2 pkg create pkg --build-type ament\_python \\
编译工作空间 & colcon build \\
加载环境 & source install/setup.bash \\
录包 & ros2 bag record -a \\
回放 bag & ros2 bag play bag \\
\bottomrule
\end{longtable}

\section{考试常考问答}

\subsection{Linux 为什么说一切皆文件？}

因为 Linux 将普通文件、目录、设备、管道、套接字等都抽象成文件，通过统一的文件接口进行访问。例如串口设备可以通过 \texttt{/dev/ttyS0} 访问，磁盘可以通过 \texttt{/dev/sda} 访问。

\subsection{绝对路径和相对路径区别是什么？}

绝对路径从根目录 \texttt{/} 开始，例如：

\begin{lstlisting}
/home/user/test.txt
\end{lstlisting}

相对路径相对于当前目录，例如：

\begin{lstlisting}
./test.txt
../src/main.c
\end{lstlisting}

\subsection{硬链接和软链接区别是什么？}

软链接类似快捷方式：

\begin{lstlisting}[language=bash]
ln -s source link_name
\end{lstlisting}

硬链接：

\begin{lstlisting}[language=bash]
ln source link_name
\end{lstlisting}

区别：
\begin{itemize}
    \item 软链接有自己的 inode，指向原路径；
    \item 硬链接和原文件共享 inode；
    \item 软链接可以跨文件系统；
    \item 硬链接通常不能链接目录；
    \item 删除原文件后，软链接失效，硬链接仍可访问数据。
\end{itemize}

\subsection{进程和线程区别是什么？}

\begin{itemize}
    \item 进程是资源分配的基本单位；
    \item 线程是 CPU 调度的基本单位；
    \item 一个进程可以包含多个线程；
    \item 同一进程内的线程共享地址空间；
    \item 进程之间相互独立，通信成本更高。
\end{itemize}

\subsection{ROS 2 Topic、Service、Action 区别是什么？}

\begin{longtable}{p{3cm}p{4cm}p{7cm}}
\toprule
类型 & 通信方式 & 适用场景 \\
\midrule
Topic & 发布订阅 & 连续数据流，如图像、雷达、速度 \\
Service & 请求响应 & 短时间任务，如查询、开关、重置 \\
Action & 目标-反馈-结果 & 长时间任务，如导航、机械臂运动 \\
\bottomrule
\end{longtable}

\subsection{ROS 2 为什么不需要 ROS Master？}

ROS 2 使用 DDS 作为底层通信机制，节点之间可以自动发现彼此，不再需要 ROS 1 中的 ROS Master 进行集中式管理。

\section{最后复习建议}

\begin{enumerate}
    \item Linux 命令一定要动手敲，尤其是 \texttt{ls、cd、cp、mv、rm、grep、find、ps、kill、chmod、tar、ssh、scp}。
    \item Shell 脚本重点掌握变量、if、for、while、函数、参数、退出码。
    \item 嵌入式 Linux 重点理解启动流程、设备文件、proc、sys、设备树、内核模块、交叉编译。
    \item ROS 2 重点理解 Node、Topic、Service、Action、Parameter、Launch、QoS。
    \item 考试遇到命令题时，先想清楚“对象是什么、操作是什么、选项是什么”。
\end{enumerate}

\end{document}